% ============================================================================
\clearpage
\section*{v4.9.0-r86 Source-Only: Independent Checks with a Second System: the Big-Bang Admixture and the Field Equations of W1}
\addcontentsline{toc}{section}{v4.9.0-r86 Source-Only: Independent Checks with a Second System: the Big-Bang Admixture and the Field Equations of W1}
\label{sec:v490-r86-wolfram-checks}
% ============================================================================

\paragraph{Purpose.}
Layers v4.9.0-r81 to r85 rest on Python and sympy scripts written in the same sessions (\S\ref{sec:v490-r81-joint-action-w1},
\S\ref{sec:v490-r85-masses-cpt}).  On 5 and 6 November 2026 a second, independent system (Wolfram Language through a connector) was used to
recompute the big-bang transition of r85.6 and to derive the field equations of the joint action W1 from their definitions.  This layer records
three pre-registered tests (H87--H89, all not blind, since the expected values were known), the errors found in the checking code, and a
correction to the reading of the energy--momentum tensor in the W1 notes.  The notes are
\texttt{notes/\allowbreak research/\allowbreak v4\_10\_0\_r86/}.  Earlier layers are not rewritten.

\noindent\textbf{Source-only status.}
This v4.9.0-r86 layer is source-only.  It does not build a PDF, does not assign a new DOI,
and does not authorize a public release.  The v4.9.0 Version DOI
\href{https://doi.org/10.5281/zenodo.23534954}{10.5281/zenodo.23534954} and the Concept DOI
\href{https://doi.org/10.5281/zenodo.19210728}{10.5281/zenodo.19210728} are unchanged by it.

\subsection*{v4.9.0-r86 results}
\label{subsec:v490-r86-results}

\paragraph{Theoremlet v4.9.0-r86.1 (H87: the admixture at the big bang, second system).}
Pre-registered (commit \texttt{61b7132}, not blind).
\begin{enumerate}[leftmargin=2em]
\item[(a)] Confirmed.  Direct integration in Wolfram (working precision 30) reproduces the values of H86 (r85.6(c)) with relative deviations
\(2.6\times10^{-13}\), \(7\times10^{-11}\), \(5\times10^{-9}\), \(5\times10^{-7}\) at \(k\Delta=2,4,6,8\); the least accurate value is the scipy value at \(k\Delta=8\).
\item[(b)] Refuted as formulated.  For large \(k\Delta\) the mode equation was integrated along \(\tau=x+ih\) below the singularity of
\(\operatorname{sech}^2\) at \(i\pi\Delta/2\) (\(h=0.9\) and \(0.95\) times \(\pi\Delta/2\)); the boundary values at \(\operatorname{Re}\tau\to\pm\infty\) are unchanged and the
problem is well conditioned (checks at \(k\Delta=4,8\) against the real axis to \(3\times10^{-10}\); both lines agree to \(1.5\times10^{-8}\)).  From
\(|C|=4.23\times10^{-12}\) at \(k\Delta=10\) to \(1.30\times10^{-37}\) at \(30\) the local slopes are \(2.925\), \(2.912\), \(2.936\), \(2.975\); between \(20\) and \(30\) the
slope is \(2.956\), \(5.9\%\) below \(\pi\), outside the \(3\%\) criterion.
\item[(c)] Confirmed.  The sharp transition of r85.5 symbolically: \(C_0=(a+b)/(2a)\), \(D_0=(a-b)/(2a)\), \(\alpha=(1+x)/(2\sqrt x)\),
\(\beta=(1-x)/(2\sqrt x)\), \(|\alpha|^2-|\beta|^2=1\); numerically \(\beta/\alpha=0.53634\), \(n=0.40382\).
\end{enumerate}
\emph{Noticed afterwards (not pre-registered).}  The local slope rises slowly on \(k\Delta\in[15,30]\).  A fit
\(\ln|C|=c_0+p\ln(k\Delta)-s\,k\Delta\) gives \(s=3.150\), \(p\approx4.8\) (exact through \(20,25,30\)) or \(s=3.077\), \(p\approx3.0\) (least squares on
\(15\)--\(30\)): compatible with \(e^{-\pi k\Delta}\) times a large power prefactor.  The position of the singularity fixes the exponent; the prefactor
depends on the strength \(\ln(a/b)-i\pi\) and is not perturbative (the first Born approximation is off by a factor \(2.2\) at \(k\Delta=8\)).  A
confirmation needs a new pre-registered test with a fixed fit model.

\paragraph{Theoremlet v4.9.0-r86.2 (H88: the field equations of W1 from their definitions).}
Pre-registered (commit \texttt{06a8380}, not blind).  All four parts confirmed, no deviation.
\begin{enumerate}[leftmargin=2em]
\item[(a)] Generator: with \(\Sigma^{ab}=\tfrac14(\sigma^a\bar\sigma^b-\sigma^b\bar\sigma^a)\), \(\Sigma^{ab\dagger}\bar\sigma^c+\bar\sigma^c\Sigma^{ab}=-(\eta^{bc}\bar\sigma^a-\eta^{ac}\bar\sigma^b)\)
in all \(64\) components, hence \(D_\mu j^c=\partial_\mu j^c+\omega_\mu{}^c{}_b\,j^b\).  The alternative ordering \(\tfrac14(\bar\sigma^a\sigma^b-\bar\sigma^b\sigma^a)\) is not
covariant.
\item[(b)] Torsion: from \(R^{ab}{}_{\mu\nu}=\partial\omega-\partial\omega+\omega\omega-\omega\omega\) and the symmetric spinor coupling, eliminating all \(24\)
contorsion components for two Weyl spinors gives exactly \(\mathcal L_{\rm eff}=-\tfrac3{16}\kappa\,j\cdot j\) with the total current (r81.1(b)).  For a
single Weyl spinor \(\mathcal L_{\rm eff}\) vanishes, since its current is null.
\item[(c)] Weyl equation: the variation with respect to \(\psi^\dagger\) is exactly
\(i\bar\sigma^\mu\partial_\mu\psi-qJA_\mu\bar\sigma^\mu\psi-\tfrac{3\kappa}{8}j_a\bar\sigma^a\psi=0\).
\item[(d)] Ward identity: \(M=qJA_\mu\bar\sigma^\mu+\tfrac{3\kappa}{8}j_a\bar\sigma^a\) is hermitian, hence \(\partial_\mu j^\mu=0\) on solutions.
\end{enumerate}

\paragraph{Theoremlet v4.9.0-r86.3 (H89: the energy--momentum tensor of W1).}
Pre-registered (commit \texttt{6355d15}, not blind).  All three parts confirmed.
\begin{enumerate}[leftmargin=2em]
\item[(a)] The vierbein variation of \(e\,\mathcal L_{4f}\) at fixed spinors gives \(T^{(4f)}_{\mu\nu}=-g_{\mu\nu}\mathcal L_{4f}=\tfrac{3\kappa}{16}\,j^2\,g_{\mu\nu}\): on
its own the four-fermion part has the form of a \(\Lambda\) term (\(p=-\rho\)), with a height \(\propto j^2\) that is not constant.
\item[(b)] Two Weyl fields of opposite spin, \(\psi_s=e^{-i(3\kappa n/8)t}\chi_s\), solve the equation of r86.2(c) exactly with total current
\((n,0,0,0)\).  On this solution \(\mathcal L_{\rm kin}=3\kappa n^2/8\), \(\mathcal L=3\kappa n^2/16\), and the whole tensor
\(T_{\mu\nu}=T^{\rm Bel}_{\mu\nu}-g_{\mu\nu}\mathcal L\) is \(\rho=p=\tfrac{3\kappa}{16}n^2>0\): stiff matter.
\item[(c)] The vierbein variation to first order around flat space, with the spin connection solved from \(de+\omega\wedge e=0\) rather than
inserted, is on a solution (two plane waves) exactly minus the Belinfante tensor (deviations \(7\times10^{-15}\) and \(5\times10^{-15}\) at values of size
\(30\)--\(90\)) and exactly symmetric.  This justifies the form used in (b).
\end{enumerate}
\emph{Disclosure.}  The first run of (c) exceeded the time limit of the connector and was reduced (spinors inserted as a solution, \(h\),
\(\partial h\) kept as symbols; no change of content).  The second run had an error in the checking code (the conjugate transpose also
conjugated the real symbols \(dh\)).  The third run disagreed (non-zero antisymmetric part); following the pre-registered order the
conventions were checked first, and the contraction \(\omega_{\mu ab}\Sigma^{ab}\) had carried extra \(\eta\) factors.  After the correction everything
agrees.  The disagreement lay in the checking code, not in W1.

\paragraph{Correction v4.9.0-r86.4 (reading of the W1 notes).}
The W1 notes of r81 (\texttt{llm/\allowbreak Wirkung\_W1\_Dixon\_U1.md}) read the four-fermion influence as ``\(p=\rho\), \(\rho\propto-\kappa n^2\), sign
convention-dependent, not like \(\Lambda\)''.  After r86.3: \(p=\rho\) is right for the whole tensor; the sign of \(\rho\) depends on the state, not on a
convention, and is positive for modes of positive frequency; ``not like \(\Lambda\)'' holds for the whole tensor, not for the four-fermion term on
its own, which cannot be separated physically from the kinetic part.  r81.1(c) (W1 produces no vacuum term) is not affected: the homogeneous
state is stiff matter, not \(\Lambda\).

\paragraph{Proof and measurement sources.}
Wolfram code \texttt{k202\_\allowbreak wolfram\_\allowbreak gegenprobe.wl}, \texttt{k203\_\allowbreak w1\_\allowbreak feldgleichungen.wl},
\texttt{k204\_\allowbreak w1\_\allowbreak einstein.wl} and the notes in \texttt{notes/\allowbreak research/\allowbreak v4\_10\_0\_r86/}; pre-registrations and
results H87--H89 in \texttt{notes/\allowbreak research/\allowbreak predictions/}.  \(\square\)

\paragraph{Consequence.}
The numbers of r85.6 are confirmed by a second system; the admixture at the big bang falls off with a rate close to \(\pi\) per unit
\(k\Delta\), with a power prefactor, so the cut-off scale is about \(1/(\pi\Delta)\) and \(\Delta\) stays free.  The field equations of W1 are derived
independently up to the full non-linear Einstein equation and the Bianchi identity: generator, torsion, Weyl equation, Ward identity, and
the energy--momentum tensor as the source to first order.  The foundation of r81--r85, as far as it rests on these equations, is firmer.
OP5 is not closed.

\subsection*{v4.9.0-r86 status in the three frames}
\label{subsec:v490-r86-frames}

\begin{itemize}[leftmargin=2em]
\item r86.1 checks the transition of r85.6 under G2 (frame T before, frame M after the big bang) with the weights of frame M; the photon frame P
gives no admixture (\(n=0\)), unchanged.
\item r86.2 and r86.3 concern the matter part of W1 in frame M.  The tachyon image under the second reading (r84) multiplies the action by the
weights \(a\), \(b\) and does not change the form of the equations; no frame-free statement is made beyond that.
\end{itemize}

\subsection*{v4.9.0-r86 register entries}
\label{subsec:v490-r86-register}

\begin{center}
\small
\begin{tabular}{@{}>{\raggedright\arraybackslash}p{0.34\textwidth}>{\raggedright\arraybackslash}p{0.58\textwidth}@{}}
\toprule
statement & status after v4.9.0-r86\\
\midrule
H86 values (r85.6(c)) & confirmed by a second system (H87a)\\
decay of the admixture & rate near \(\pi\) with power prefactor (afterwards); \(3\%\) criterion missed (H87b)\\
W1 generator, torsion, Weyl, Ward & derived from the definitions (H88)\\
energy--momentum tensor of W1 & minus Belinfante to first order; homogeneous state \(\rho=p>0\) (H89)\\
four-fermion term alone & \(\Lambda\) form with non-constant height; W1 vacuum term still absent\\
full Einstein equation, Bianchi & open\\
\bottomrule
\end{tabular}
\end{center}

\subsection*{v4.9.0-r86 ledger and guard update}
\label{subsec:v490-r86-ledger-guard}

\[
R_{\rm chk}^{v490\text{-}r86}=(1,1,0)
\]
for (r85.6 confirmed independently; W1 field equations derived to first order; full Einstein equation open).

The layer spends no guard:
\[
R_{\rm glob}^{(v490\text{-}r86,+)}=(0,1,1,1,1),
\qquad
\Gamma_{\rm guard}^{v490\text{-}r86}=(1,1,1,1,1),
\qquad
\beta_4=1/256\ \text{unspent}.
\]
It does not imply exact finite \(C5\) zero closure, exact all-readout terminality,
\(C6/T2\), a global photon-clock/RH theorem, or global all-readout closure.
